\chapter{FX Derivatives}
\label{ch:fx_derivatives}
\chaptermark{FX derivatives}

Foreign-exchange (FX) is the largest financial market in the world:
the BIS 2022 Triennial Survey put average daily turnover at \$7.5
trillion, an order of magnitude above global equities.  It also
carries the densest jargon --- ``pips'', ``forward points'',
``25-delta risk-reversal'', ``TARF'', ``CNH vs CNY'' --- most with
no analogue elsewhere.  This chapter translates the jargon, derives
\emph{covered interest parity}, states the FX Black--Scholes
(\emph{Garman--Kohlhagen}), and tours the vanilla-through-exotic
option zoo: barriers, digitals, TARFs, and the ATMF/RR/BF quoting
convention from which the entire smile is reconstructed.

\begin{backgroundbox}[Prerequisites]
From earlier in the book:
\begin{itemize}
  \item \Cref{ch:arbitrage_zoo}: the no-arbitrage principle, triangular
    FX arbitrage, and the idea that cross-rates must agree in a
    cycle.
  \item \Cref{ch:options_payoffs}: European call/put payoffs,
    put--call parity, and the notion of a strike $\Kstrike$.
  \item \Cref{ch:black_scholes}: the Black--Scholes formula under
    geometric Brownian motion, the $\CDFN$/$\pdfN$ notation, and
    risk-neutral pricing.
  \item \Cref{ch:fixed_income_foundations}: continuous discounting,
    zero-coupon curves, $e^{-\rrf T}$ as the discount factor.
\end{itemize}
If any of these feel hazy, re-read the relevant chapter first.
\end{backgroundbox}


%% ================================================================
\section{FX jargon decoder}
\label{sec:fx_jargon}
%% ================================================================

FX desk conversation is unreadable without the shorthand.  The
table collects the terms used in the rest of the chapter; each is
formally defined where first needed.

\begin{keyfactbox}[FX jargon at a glance]
\label{fact:fx_jargon}
\begin{tabular}{@{}p{0.22\textwidth} p{0.72\textwidth}@{}}
\toprule
\textbf{Term} & \textbf{Meaning} \\
\midrule
\emph{Currency pair} & A pair written $\mathrm{BASE}/\mathrm{QUOTE}$,
  e.g.\ $\mathrm{EUR}/\mathrm{USD}$.  The quoted number is the number
  of units of the \emph{quote} currency required to buy one unit of
  the \emph{base} currency. \\
\emph{Spot} $S_t$ & Current exchange rate; deliverable on the
  \emph{spot date}, conventionally $T+2$ business days. \\
\emph{Forward} $F(T)$ & Rate agreed now for an exchange settling at
  future date $T$.  No money changes hands today. \\
\emph{Forward points} & The difference $F(T) - S_t$, quoted in
  \emph{pips}. \\
\emph{Pip} & ``Percentage in point''; the smallest conventional
  price increment.  For most pairs this is $10^{-4}$; for JPY pairs it
  is $10^{-2}$. \\
\emph{Direct quote} & Quote with the domestic currency as the quote
  (right-hand) currency.  From a US perspective, $\mathrm{EUR}/\mathrm{USD}$
  is direct; $\mathrm{USD}/\mathrm{JPY}$ is indirect. \\
\emph{Cross rate} & Rate for a pair neither leg of which is the USD,
  e.g.\ $\mathrm{EUR}/\mathrm{JPY}$, usually inferred from the two
  USD legs. \\
\emph{Domestic / foreign} & In option pricing, the \emph{domestic}
  currency is the one you are measuring P\&L in (the quote currency);
  the \emph{foreign} currency is the underlying asset (the base
  currency).  Rates $r_d, r_f$ follow this convention. \\
\emph{NDF} & Non-deliverable forward: cash-settles in USD against a
  fixing, used where the local currency is not freely convertible
  (e.g.\ $\mathrm{KRW}, \mathrm{TWD}, \mathrm{CNY}$). \\
\emph{CNH vs CNY} & Offshore (Hong Kong) vs onshore (mainland)
  renminbi.  The two are separate markets and their rates can
  diverge persistently. \\
\bottomrule
\end{tabular}
\end{keyfactbox}

\begin{pitfallbox}[$\mathrm{EUR}/\mathrm{USD}$ is not the same as $\mathrm{USD}/\mathrm{EUR}$]
$\mathrm{EUR}/\mathrm{USD} = 1.08$ means
``\textbf{1} EUR buys \textbf{1.08} USD.''  Inverting, one USD buys
$1/1.08 \approx 0.9259$ EUR.  Newcomers flip the convention and
wonder why their hedge is backwards.  The rule:
\emph{the left-hand currency is the one unit; the right-hand is
how many you pay.}  ``A call on $\mathrm{EUR}/\mathrm{USD}$'' is a
call on EUR --- buying EUR, selling USD --- struck in USD per EUR.
\end{pitfallbox}

The slot convention is fixed by market practice, not mathematics.
EUR is always the base against USD, JPY is always the quote, GBP is
always the base, and so on.  Get the convention wrong on a ticket
and you have written the opposite trade.

\subsection{Fixings and settlement mechanics}
\label{subsec:fx_fixings}

Many FX contracts --- NDFs, option barriers, TARF schedules,
hedged-equity baskets --- reference an \emph{FX fixing}: a
publicly-observed rate set at a fixed time by polling or matching
on a specified venue.  The main fixings:
\begin{itemize}
  \item \textbf{WM/Reuters 4pm London}: benchmark for equity-index
    currency hedging; set from trades and executable quotes in a
    five-minute window around 4pm London, published by Refinitiv.
  \item \textbf{ECB reference rates}: daily concertation fix at
    14:15 CET, for European statistics and accounting.
  \item \textbf{Central-bank NDF fixes}: e.g.\ the PBOC USD/CNY
    onshore midpoint (09:15 Beijing), reference for many CNY NDFs.
\end{itemize}
Confusing which fixing a contract references has cost traders real
money: a barrier that knocks out ``on any touch'' does so
continuously on spot, but a TARF leg ``settling against the fixing''
cares only about the one reference value on the fixing date.

\textbf{Non-deliverable forwards (NDFs).}
An NDF is a forward with no physical exchange at settlement; one
party pays the other USD (or EUR) cash equal to (forward rate --
fixing) $\times$ notional.  NDFs exist because many EM currencies
(CNY, KRW, INR, TWD, BRL, \ldots) have capital controls preventing
free offshore ownership of the local currency; the NDF gives
exposure without holding the local leg.  The NDF market is the
primary venue for non-resident renminbi, rupee, and won exposure,
with over \$250bn daily turnover in 2022.


%% ================================================================
\section{Spot and forwards}
\label{sec:fx_spot_forward}
%% ================================================================

\subsection{Spot}

The \textbf{spot rate} $S_t$ is the price at time $t$ for immediate
exchange of the two currencies.  ``Immediate'' is convention: for
most pairs the \emph{spot date} is $T+2$ business days.  USD/CAD is
$T+1$; intraday fixings use whatever the fixing definition says.

The $T+2$ convention is historical: pre-electronic settlement
required physical payment instructions through two central-bank
RTGS systems in two time zones.  CLS eliminates the operational
need, but $T+2$ remains the anchor --- zero curves, forward-point
quotes, and option-expiry calendars all hang off it.  A ``tom/next''
trade is a one-day forward relative to spot, not a two-day one.

\subsection{Forwards}

A \textbf{forward contract} struck at $t$ exchanges currencies at
fixed rate $F(t,T)$ on date $T$.  No money changes hands at
inception; at $T$ one side delivers foreign and receives domestic at
$F$.  The commercial motivation is hedging: a European exporter
invoiced USD 10M for delivery in three months faces uncertain EUR
proceeds, and locking a forward today converts that uncertain
future USD receipt into a known EUR amount.

What rate makes this fair today?  No-arbitrage pins it down exactly
as with equity forwards (\cref{ch:options_payoffs}), with one twist:
holding the foreign currency earns the \emph{foreign interest rate}
$r_f$, not a dividend yield.

\subsection{Covered interest parity}
\label{subsec:cip}

\begin{keyfactbox}[Covered interest parity (CIP)]
\label{fact:cip}
Let $S_t$ be the spot $\mathrm{BASE}/\mathrm{QUOTE}$ rate (units of
quote per unit of base).  Let $r_d$ be the continuously compounded
domestic (quote-currency) interest rate and $r_f$ the foreign
(base-currency) rate, both for the tenor $T$.  Then the fair forward
rate is
\begin{equation}
  F(t,T) \;=\; S_t \, e^{(r_d - r_f)(T - t)}.
  \label{eq:cip}
\end{equation}
If market quotes deviate from this relation, a textbook arbitrage
exists.
\end{keyfactbox}

\textbf{Derivation.}  Consider two self-financing strategies,
starting from one unit of foreign at time $t$ and ending in domestic
at $T$.

\textit{Strategy A (invest foreign, convert forward):} deposit the
foreign unit at rate $r_f$, growing to $e^{r_f (T-t)}$; sell this
forward at $F(t,T)$; receive $F(t,T)\,e^{r_f(T-t)}$ domestic at $T$.

\textit{Strategy B (convert spot, invest domestic):} sell the
foreign unit at spot for $S_t$ domestic; deposit at $r_d$, growing
to $S_t \, e^{r_d(T-t)}$ at $T$.

Both start from one foreign unit with deterministic terminal
cashflow.  By the law of one price they must agree:
\begin{equation}
  F(t,T) \, e^{r_f(T-t)} \;=\; S_t \, e^{r_d(T-t)},
\end{equation}
rearranging to \cref{eq:cip}.  If $F$ exceeded this level, borrow
domestic, run Strategy B, sell forward, and collect
$Fe^{r_f(T-t)} - S_t e^{r_d(T-t)} > 0$ risklessly; the converse
flips.  \citeHull, Ch.\ 5 (general assets) and Ch.\ 17 (currencies).

\begin{intuitionbox}
$F > S$ iff the domestic rate exceeds the foreign rate.  If USD
earns 5\% vs EUR's 4\%, everyone wants USD; the forward market
compensates by pricing future USD \emph{lower} in EUR terms so the
two investments break even.  Mnemonic: \emph{the high-yielding
currency trades at a forward discount.}
\end{intuitionbox}

\begin{figure}[htbp]
\centering
\begin{tikzpicture}[>=Stealth,
  box/.style={draw, rounded corners, fill=gray!8,
              minimum width=3.2cm, minimum height=1.1cm,
              align=center, font=\footnotesize}]
  \node[box] (A) at (0,   2.5) {1 unit foreign\\at time $t$};
  \node[box] (B) at (7.5, 2.5) {$e^{r_f(T-t)}$ foreign\\at time $T$};
  \node[box] (C) at (0,   0)   {$S_t$ domestic\\at time $t$};
  \node[box] (D) at (7.5, 0)   {$S_t\,e^{r_d(T-t)}$\\domestic at $T$};

  % Strategy A: invest foreign, sell forward
  \draw[->, thick, blue] (A) -- node[above, font=\footnotesize]{invest at $r_f$} (B);
  \draw[->, thick, blue] (B) -- node[right, align=center, font=\footnotesize]
    {sell forward\\at $F(t,T)$} (D);

  % Strategy B: sell spot, invest domestic
  \draw[->, thick, red] (A) -- node[left, font=\footnotesize]{sell spot at $S_t$} (C);
  \draw[->, thick, red] (C) -- node[below, font=\footnotesize]{invest at $r_d$} (D);

  % Strategy labels
  \node[blue, font=\footnotesize\bfseries] at (3.75, 3.15) {Strategy A};
  \node[red,  font=\footnotesize\bfseries] at (3.75,-0.65) {Strategy B};

  % No-arb result
  \node[font=\footnotesize, align=center] at (3.75, 1.25)
    {No-arb: equal terminal values\\
     $\Rightarrow F(t,T) = S_t\,e^{(r_d - r_f)(T-t)}$};
\end{tikzpicture}
\caption{Covered interest parity (CIP). Both strategies start from 1 unit of foreign
currency at $t$ and end in domestic at $T$. Strategy A (blue): deposit foreign at $r_f$,
sell proceeds forward at $F$. Strategy B (red): sell spot at $S_t$, invest domestic at
$r_d$. No-arbitrage forces the two terminal values to agree, pinning the forward rate.}
\label{fig:cip_strategies}
\end{figure}

\subsection{Worked example: 1-year EUR/USD forward}
\label{subsec:cip_example}

Suppose EUR/USD spot $S_0 = 1.0800$, the continuously compounded USD
rate is $r_d = 5\%$, and the continuously compounded EUR rate is
$r_f = 4\%$, both for 1-year tenor.  Then
\begin{equation}
  F(0, 1) \;=\; 1.08 \cdot e^{(0.05 - 0.04) \cdot 1}
            \;=\; 1.08 \cdot e^{0.01}
            \;\approx\; 1.0909.
\end{equation}
The forward points are $F - S = 0.010854 = 108.54$ pips, quoted as
``EUR/USD 1-year fwd at $+108.5$ points.''  Verified to six decimal
places in Python (\cref{lst:gk_pricer}).

\begin{pitfallbox}[Quote convention flip reverses the sign]
With $\mathrm{EUR}/\mathrm{USD}$ and USD domestic, higher USD rate
gives a forward \emph{premium} ($F > S$).  Quoted as
$\mathrm{USD}/\mathrm{EUR}$ instead, the domestic role swaps:
$r_d = 4\%, r_f = 5\%$, so $F/S = e^{-0.01}$ --- a forward
\emph{discount}.  Same economics, opposite sign on the points.
\end{pitfallbox}


%% ================================================================
\section{Cross-currency basis swaps}
\label{sec:xccy_basis}
%% ================================================================

A \textbf{cross-currency basis swap} (XCCY swap) exchanges floating
interest in two currencies over months to decades.  At inception,
the parties exchange principals $N$ and $S_0 N$ at spot $S_0$ (one
pays $N$ USD, receives $S_0 N$ EUR).  At each coupon $t_i$, one pays
floating USD (SOFR) on $N$ and receives floating EUR (EURIBOR plus
spread) on $S_0 N$, spread on the EUR leg.  At maturity, principals
are re-exchanged at the \emph{initial} rate $S_0$, not the
prevailing one.

The ``basis'' $b$ is the spread on the non-USD leg that zeros
inception PV.  Under naive CIP this would be zero; since 2008 it has
been persistently non-zero, typically 10--50~bp against USD for
majors and wider in stress.  In plain English: $b$ is the extra
interest a counterparty demands, on top of their own domestic
floating rate, to lend you their currency and take USD in exchange
--- compensation for the scarcity of dollar funding when non-US
banks want USD more than US banks want foreign currency.

\begin{keyfactbox}[Post-2008 CIP deviations]
\label{fact:cip_deviations}
The observed cross-currency basis $b$ violates pure CIP:
\emph{borrowing USD synthetically via FX swaps is more expensive
than borrowing USD directly.}  The drivers are
\begin{itemize}
  \item \textbf{Balance-sheet costs} under Basel III --- arbitraging
    the basis consumes scarce leverage capacity.
  \item \textbf{USD funding premium} --- non-US banks structurally
    short of USD bid for it, especially at quarter- and year-ends.
  \item \textbf{Credit/counterparty risk} in the FX-swap market.
\end{itemize}
See \citet{borio2016covered} for the canonical post-crisis account.
For option-pricing purposes, the practical consequence is that the
forward rate used in \cref{fact:cip} must be replaced with the
\emph{market} forward (which bakes in the basis), not the
textbook CIP-implied forward.
\end{keyfactbox}

The basis is a term structure quoted at 3M, 6M, 1Y, 2Y, \ldots, 30Y.
Long-tenor FX-forward desks price off the basis curve plus the two
single-currency curves, not textbook CIP.  XCCY basis swaps hedge
and express views on the spread.

\subsection{Implied basis from a forward}
\label{subsec:implied_basis}

Given market forward $F^{\mathrm{mkt}}$ and two zero curves, the
\emph{implied} basis is
\begin{equation}
  b_{\mathrm{implied}}(T) \;=\;
  \frac{1}{T}\ln\frac{F^{\mathrm{mkt}}}{S_0}
  \;-\; (r_d - r_f),
  \label{eq:implied_basis}
\end{equation}
the deviation of the observed forward from the textbook CIP level
$S_0 e^{(r_d-r_f)T}$.  Positive = premium over CIP; negative =
discount.  EUR/USD 3-month basis has blown out to 50--80~bp at
stressed quarter-ends.

\subsection{Why this matters for pricing}

The Garman--Kohlhagen formula of \cref{sec:fx_vanillas} assumes CIP
and uses $r_d - r_f$ as the risk-neutral drift.  Desks pricing
against market forwards replace this with
$\tfrac{1}{T}\ln(F^{\mathrm{mkt}}/S_0) = r_d - r_f +
b_{\mathrm{implied}}$.  Equivalently, price off the market forward
curve directly (the \emph{forward-flat} convention).


%% ================================================================
\section{FX vanillas and quanto}
\label{sec:fx_vanillas}
%% ================================================================

\subsection{Garman--Kohlhagen}

The FX analogue of Black--Scholes is \textbf{Garman--Kohlhagen}
\citep{garman1983foreign}: treat the foreign currency like a
dividend-paying stock with continuous yield $r_f$ in place of $q$.

\begin{keyfactbox}[Garman--Kohlhagen formula]
\label{fact:garman_kohlhagen}
Under the risk-neutral model
$d\Spx_t = (r_d - r_f)\,\Spx_t\,dt + \volat\,\Spx_t\,dW_t$,
a European call on the FX rate $\Spx$ (paying
$\max(\Spx_T - \Kstrike, 0)$ in domestic currency) is worth
\begin{equation}
  \Call \;=\; \Spx_0 \, e^{-r_f T} \CDFN(d_1)
              \;-\; \Kstrike \, e^{-r_d T} \CDFN(d_2),
  \label{eq:gk_call}
\end{equation}
with
\begin{equation}
  d_1 \;=\; \frac{\ln(\Spx_0/\Kstrike) + (r_d - r_f + \tfrac12 \volat^2)T}
                 {\volat \sqrt{T}},
  \qquad
  d_2 \;=\; d_1 - \volat\sqrt{T}.
  \label{eq:gk_d1d2}
\end{equation}
The corresponding put price is
$\Put = \Kstrike e^{-r_d T}\CDFN(-d_2) - \Spx_0 e^{-r_f T}\CDFN(-d_1)$.
\end{keyfactbox}

Differences from equity Black--Scholes (\cref{ch:black_scholes}) are
cosmetic: drift under $\Qm$ is $r_d - r_f$; spot discounted by
$e^{-r_f T}$.  The derivation is identical with
$q \mapsto r_f$ throughout.  \citeHull, Ch.\ 17.

\subsection{Put--call parity in FX form}

Subtracting put from call gives FX put--call parity
\begin{equation}
  \Call - \Put \;=\; \Spx_0 e^{-r_f T} - \Kstrike e^{-r_d T},
  \label{eq:fx_putcall}
\end{equation}
the PV of a forward buying foreign at $\Kstrike$.  Using CIP
(\cref{fact:cip}), the RHS equals $e^{-r_d T}(F - \Kstrike)$.  Long
call + short put at the same strike = long forward, as in equities.

\subsection{Worked example: 3-month EUR/USD ATM call}
\label{subsec:gk_example}

Using the same market data as in \cref{subsec:cip_example} but with
$T = 0.25$ years, $\volat = 9\%$ (a typical EUR/USD at-the-money
implied vol), and strike $\Kstrike = 1.08$ (ATM spot):
\begin{align}
  d_1 &= \frac{\ln(1) + (0.05 - 0.04 + \tfrac12 \cdot 0.09^2) \cdot 0.25}
               {0.09 \sqrt{0.25}}
        \;=\; 0.0781, \\
  d_2 &= 0.0781 - 0.09\sqrt{0.25} \;=\; 0.0331, \\
  \Call &= 1.08\,e^{-0.04\cdot 0.25}\CDFN(0.0781)
           - 1.08\,e^{-0.05\cdot 0.25}\CDFN(0.0331) \notag \\
        &\approx 0.02054 \text{ USD per EUR notional}.
\end{align}
On a EUR 10M notional, the premium is USD 205{,}346.
The corresponding put is worth USD 178{,}648; their difference
USD 26{,}698 matches $\Spx e^{-r_f T} - \Kstrike e^{-r_d T}$ times
EUR 10M, as parity requires.  All five numbers above are
verified in Python below to six decimal places.

\begin{lstlisting}[caption={Garman--Kohlhagen vanilla pricer for FX options.},label={lst:gk_pricer}]
import numpy as np
from scipy.stats import norm

def garman_kohlhagen(S, K, T, r_d, r_f, sigma, option="call"):
    """Price a European FX option (domestic-currency premium per
    unit of foreign notional).  S, K in domestic-per-foreign units."""
    d1 = (np.log(S / K) + (r_d - r_f + 0.5 * sigma**2) * T) \
         / (sigma * np.sqrt(T))
    d2 = d1 - sigma * np.sqrt(T)
    if option == "call":
        return S * np.exp(-r_f * T) * norm.cdf(d1) \
             - K * np.exp(-r_d * T) * norm.cdf(d2)
    return K * np.exp(-r_d * T) * norm.cdf(-d2) \
         - S * np.exp(-r_f * T) * norm.cdf(-d1)

# 3M EUR/USD ATM, spot 1.08, vol 9%, USD 5%, EUR 4%
C = garman_kohlhagen(1.08, 1.08, 0.25, 0.05, 0.04, 0.09, "call")
P = garman_kohlhagen(1.08, 1.08, 0.25, 0.05, 0.04, 0.09, "put")
print(f"Call = {C:.6f}, Put = {P:.6f}, C-P = {C-P:.6f}")
# Check put-call parity: C - P == S*exp(-r_f*T) - K*exp(-r_d*T)
parity = 1.08*np.exp(-0.04*0.25) - 1.08*np.exp(-0.05*0.25)
assert abs((C - P) - parity) < 1e-10
\end{lstlisting}

\subsection{Quanto options}

A \textbf{quanto} (``quantity-adjusting'') option pays in a currency
different from the underlying's natural one.  Canonical example: a
Japanese investor holds a call on SPX paying
$\max(\Spx_T^{\mathrm{SPX}} - \Kstrike, 0)$ \emph{in JPY}, with a
contractually fixed FX rate (say 1 index unit = 1 JPY).

Quantos introduce correlation between index $\Spx^{\mathrm{SPX}}$
and rate $X^{\mathrm{USDJPY}}$.  Under the JPY risk-neutral measure,
the SPX drift is
\begin{equation}
  \mu^{\mathrm{SPX, quanto}} \;=\;
  r_{\mathrm{JPY}} - q^{\mathrm{SPX}}
  \;-\; \rho \, \volat_{\mathrm{SPX}} \, \volat_{\mathrm{USDJPY}},
  \label{eq:quanto_drift}
\end{equation}
with $\rho$ the correlation of SPX and USD/JPY log-returns.  The
adjustment is the $-\rho\volat\volat$ term: positive co-movement
reduces the quanto drift, because a JPY investor gets ``free'' FX
appreciation when SPX rises and pays for it via a lower forward.
With the adjusted drift the rest is ordinary Black--Scholes.
\citeHull, Ch.\ 30.

\subsubsection{Worked example: quanto vs naive SPX call}

Take $\volat_{\mathrm{SPX}} = 20\%,
\volat_{\mathrm{USDJPY}} = 10\%, \rho = -0.30$ (USD weakens when
SPX rallies).  The SPX-drift adjustment is
\begin{equation}
  -\rho\,\volat_{\mathrm{SPX}}\,\volat_{\mathrm{USDJPY}}
  \;=\;
  -(-0.30)(0.20)(0.10) \;=\; +0.60\%,
\end{equation}
raising the quanto drift by 60 bp/year.  The JPY investor pays
\emph{more} than for the unhedged exposure because she gets the
FX-hedged return without paying for the hedge.  Flip $\rho > 0$ and
the quanto trades at a discount.

\begin{pitfallbox}[Quantos are not the same as FX-hedged notes]
An \emph{FX-hedged} investment converts the SPX return at the
terminal FX rate and strips FX exposure via a rolling forward hedge.
A \emph{quanto} uses a contractually fixed FX rate --- no spot-FX
sensitivity at expiry.  The P\&L differs, and the quanto carries
correlation risk ($\rho$) that the hedged note does not.
\end{pitfallbox}


%% ================================================================
\section{Barriers, digitals, and TARFs}
\label{sec:fx_exotics}
%% ================================================================

Vanilla FX options trade in vast quantities, but exotic-desk
complexity lives in structured products used to fine-tune corporate
and asset-manager hedges.  Three families follow.

\subsection{Barrier options}

A \textbf{barrier option} is a vanilla that comes into existence
(\emph{knock-in}) or is extinguished (\emph{knock-out}) when spot
touches barrier $B$ during the option's life.  \{up, down\} $\times$
\{in, out\} $\times$ \{call, put\} gives eight standard types, all
with closed-form GK prices
\citep[see][Ch.\ 26 and 19]{hull8ed}.  The key \emph{in-out parity}:
\begin{equation}
  \text{knock-in} + \text{knock-out} \;=\; \text{vanilla},
  \label{eq:in_out_parity}
\end{equation}
since the barrier either fires (knock-in lives, knock-out dies) or
does not, and the portfolio pays the vanilla payoff either way.

Barriers are popular because they are strictly cheaper than the
matching vanilla: a hedger who ``knows'' EUR/USD won't rally above
1.20 in six months buys a 1.10-strike knock-out with barrier 1.20.
If spot touches 1.20, the hedge disappears exactly when most needed
--- the standard barrier failure mode.

For the case of an \emph{up-and-out call} with strike $\Kstrike <
B$ under Garman--Kohlhagen dynamics, the closed form is
\citep[][Ch.\ 26]{hull8ed}
\begin{equation}
  \Call_{\mathrm{UO}} \;=\;
  \Call_{\mathrm{vanilla}}
  \;-\; \Call_{\mathrm{UI}},
  \label{eq:uo_decomp}
\end{equation}
where $\Call_{\mathrm{UI}}$ is a sum of four Gaussian CDF terms
involving the reflected strike $B^2/\Spx_0$ and exponent
$\mu = (r_d - r_f)/\volat^2 - \tfrac12$.  Ugly but mechanical to
code.  Key properties:
\begin{itemize}
  \item UO vanishes when spot hits $B$;
  \item vega changes sign near $B$ (long far from $B$, short close
    --- more vol raises the knockout probability);
  \item gamma spikes large and negative near $B$, making delta
    hedging unstable in the final 1--2\% approach.
\end{itemize}
The last point --- barrier \emph{pin risk} --- is why traders run
separate ``stopping-time'' hedges (small digital offsets for the
gamma spike) rather than smooth delta hedging alone.

\subsubsection{Barrier vs vanilla: a concrete comparison}

Using the data of \cref{subsec:gk_example}
(spot 1.08, $r_d = 5\%, r_f = 4\%, \volat = 9\%, T = 0.25$), compare
three $\Kstrike = 1.09$ calls:
\begin{itemize}
  \item \textbf{Vanilla}: $\approx 1.59\%$ of EUR notional.
  \item \textbf{Up-and-out, $B = 1.12$}: $\approx 0.11\%$ (MC).  An
    order of magnitude cheaper: the barrier is tight enough that
    most paying paths knock out first.
  \item \textbf{Up-and-out, $B = 1.15$}: $\approx 0.61\%$ ---
    under half the vanilla, but retains most of the 1.09--1.15
    terminal probability.
\end{itemize}
(Prices verified via MC at $2\times 10^5$ paths, 100 steps.)
The discount is largest exactly when the cheap hedge is most
wanted --- the same perverse trade-off as TARFs.  The barrier's
expected payoff conditional on the hedge being needed
(spot $\geq 1.09$ at expiry) is strictly lower than the vanilla's;
the shortfall is the structure's true cost.

\subsection{Digital (binary) options}

A \textbf{digital call} pays $Q$ if $\Spx_T > \Kstrike$, zero
otherwise.  Under GK, the domestic-cash digital call is
\begin{equation}
  D^{\mathrm{call}}_{\mathrm{cash}} \;=\;
  Q \, e^{-r_d T} \, \CDFN(d_2),
  \label{eq:digital_call}
\end{equation}
with $d_2$ as in \cref{eq:gk_d1d2}.  The digital put is
$Q e^{-r_d T} \CDFN(-d_2)$; they sum to $Q e^{-r_d T}$, the digital
put--call parity.

\begin{intuitionbox}
A digital is the derivative of a vanilla with respect to strike:
$\partial C / \partial \Kstrike = -e^{-r_d T}\CDFN(d_2)$, so a tiny
long-call-short-call-at-$\Kstrike + \epsilon$ spread, rescaled by
$1/\epsilon$, replicates a digital.  This is why a digital's
\emph{vega profile} has a sharp spike at the strike: the replicating
call spread becomes infinitely narrow.  Traders describe digitals as
``event insurance'' --- they pay a fixed amount contingent on a
spot-at-expiry event.
\end{intuitionbox}

``One-touch'' and ``no-touch'' variants pay $Q$ if spot ever / never
crosses the barrier before expiry --- cheap conditional bets on
range-bound or trending regimes.  Closed forms follow from the
reflection principle \citep[][Ch.\ 19]{hull8ed}.

\subsubsection{One-touch pricing sketch}

For a one-touch that pays $Q$ in domestic currency at expiry $T$
if spot ever hits barrier $B$, the BS price under drift
$\mu = r_d - r_f$ is
\begin{equation}
  \mathrm{OT} \;=\; Q\,e^{-r_d T}\!
    \left[ \left(\tfrac{B}{\Spx_0}\right)^{2\mu/\volat^2}\!
      \CDFN(z_+)
    + \CDFN(z_-) \right],
  \label{eq:one_touch}
\end{equation}
with
$z_\pm = (\pm \ln(\Spx_0/B) - \mu T + \tfrac12\volat^2 T)/(\volat\sqrt{T})$
(sign conventions vary; \citet{clark2011foreign} lists all four
variants).  One-touches are pure \emph{stopping-time} products:
price depends on the first-passage distribution of GBM, not the
terminal alone.  Skew and kurtosis drive large BS corrections, so
desks mark them via Vanna--Volga or smile-consistent MC.

\subsection{Target redemption forwards (TARFs)}

A \textbf{target redemption forward} is a structured series of
monthly or weekly forwards: on each fixing, the client buys base
currency at a favourable strike if spot is in range.  The structure
\emph{terminates early} once cumulative client gains reach a
contractual \emph{target}; if spot moves against the client, they
keep paying at the unfavourable strike, often on \emph{doubled}
notional, through the remaining schedule.

This asymmetry is the point: clients accept tail risk in exchange
for a better favourable-side strike.  Dealers price TARFs by MC over
the full strike schedule with the early-termination trigger; no
closed form.

Concretely, a 12-month EUR/USD TARF might specify: twelve monthly
fixings, strike $\Kstrike = 1.09$, notional EUR 1M per fixing,
target cumulative client P\&L = EUR 100k, leverage 2x on the loss
leg.  On each fixing date the cashflow to the client is
\begin{equation}
  \mathrm{CF}_i \;=\;
  \begin{cases}
    (\Kstrike - \mathrm{Fix}_i) \cdot N
      & \text{if } \mathrm{Fix}_i \leq \Kstrike \ \text{(client gain)}, \\
    (\Kstrike - \mathrm{Fix}_i) \cdot 2N
      & \text{if } \mathrm{Fix}_i > \Kstrike \ \text{(client loss, 2x)},
  \end{cases}
\end{equation}
with $N = $ EUR 1M.  Termination is at the first fixing $i^*$ with
$\sum_{j \leq i^*} \max(\mathrm{CF}_j, 0) \geq $ target.  The MC
pricer draws 12 fixing paths, applies the termination rule,
discounts, and averages.  Greeks come from bump-and-reprice at
roughly $12\times$ vanilla cost.

\begin{pitfallbox}[TARFs are leveraged the wrong way]
TARFs cap client upside (the target) while leveraging downside
tenor-unlimited --- the wrong profile for a hedger, the right one
for a short-vol speculator.  Treasurers were sold TARFs as ``cheap
hedges'' in 2014--2015; the August 2015 RMB devaluation locked
Asian corporates into daily-doubling CNH purchases above market,
booking multi-hundred-million USD losses.  The product is now
regulatory-flagged in most jurisdictions.  \emph{If a structure
advertises ``free'' hedging, find the short-vol double-notional
leg before signing.}
\end{pitfallbox}

\begin{historybox}[title={The 2015 CNH TARF blow-up}]
July 2015: 1Y CNH near 6.20 per USD with a mild appreciation
trend.  Asian treasurers had TARFs selling CNH at 6.15 (0.8\% better
than spot), with doubled notional if CNH depreciated.  On 11 August
the PBOC devalued the onshore fix by 1.9\%; CNH fell to 6.40 over
weeks.  Corporates sold double-notional at 6.15 into spot near 6.40
--- roughly 4\%/month MTM loss on twice notional, compounding.
Aggregate Asian corporate losses exceeded \$10bn --- the textbook
asymmetric-payoff mis-selling case.
\end{historybox}


%% ================================================================
\section{Delta conventions: ATMF, 25\texorpdfstring{$\Delta$}{-delta} RR, 25\texorpdfstring{$\Delta$}{-delta} BF}
\label{sec:fx_delta_conventions}
%% ================================================================

FX options are quoted by \emph{delta}, not strike or moneyness.  A
broker quote for EUR/USD 1-month looks like
\begin{center}
\small
\begin{tabular}{@{}lll@{}}
\toprule
1M ATMF & 8.50\% & (at-the-money-forward vol) \\
1M 25\,RR & $-0.40\%$ & (25-delta risk reversal) \\
1M 25\,BF & $+0.18\%$ & (25-delta butterfly) \\
\bottomrule
\end{tabular}
\end{center}
From these three numbers the entire 1-month smile is reconstructed.

\subsection{Delta and ATMF}

\textbf{Delta} is $\Delta = \partial V / \partial S$.  ``25-delta
call'' is the strike with $\Delta = 0.25$; ``25-delta put'' has
$\Delta = -0.25$.  \citet{clark2011foreign} gives the full taxonomy
of conventions:
\begin{itemize}
  \item \textbf{Spot delta, premium-excluded (SD-PE)}: plain
    $\partial V / \partial S$.  Default for USD-premium pairs
    (EUR/USD, GBP/USD).
  \item \textbf{Spot delta, premium-included (SD-PI)}: adjusted
    for foreign-currency premium payment.  Standard for USD/JPY etc.
  \item \textbf{Forward delta (FD-PE, FD-PI)}: $\partial V/\partial
    F$; standard for long-dated options ($T > 1$ year).
\end{itemize}
One option, four numerical deltas; desks exchanging a ``25-delta
call'' without fixing the convention can be pricing different
strikes.

The \textbf{at-the-money forward} (ATMF) strike is
$\Kstrike_{\mathrm{ATMF}} = F(0,T) = S_0 e^{(r_d - r_f)T}$, the
delta-neutral straddle strike under most conventions.  ATMF vol is
the single most-quoted number.  ``ATM-spot'' ($\Kstrike = S_0$)
differs from ATMF whenever $r_d \neq r_f$; FX defaults to ATMF.

\subsection{Risk reversal and butterfly}

\begin{keyfactbox}[FX vol quoting via ATMF/RR/BF]
\label{fact:fx_rr_bf}
For a fixed delta $\delta$ (typically $25\%$ or $10\%$), define
\begin{align}
  \mathrm{RR}_\delta
    &\;=\; \volat(\delta\text{-call}) - \volat(\delta\text{-put}),
    \label{eq:rr_def} \\
  \mathrm{BF}_\delta
    &\;=\; \tfrac12\!\left[\volat(\delta\text{-call})
                            + \volat(\delta\text{-put})\right]
           - \volat_{\mathrm{ATMF}}.
    \label{eq:bf_def}
\end{align}
$\mathrm{RR}_\delta$ measures the \emph{skew} (asymmetry of the
smile); $\mathrm{BF}_\delta$ measures the \emph{smile convexity}
(curvature).  Given $(\volat_{\mathrm{ATMF}}, \mathrm{RR}_\delta,
\mathrm{BF}_\delta)$, the wing vols are recovered by
\begin{equation}
  \volat(\delta\text{-call}) \;=\;
    \volat_{\mathrm{ATMF}} + \mathrm{BF}_\delta + \tfrac12\mathrm{RR}_\delta,
  \qquad
  \volat(\delta\text{-put}) \;=\;
    \volat_{\mathrm{ATMF}} + \mathrm{BF}_\delta - \tfrac12\mathrm{RR}_\delta.
  \label{eq:wing_vols}
\end{equation}
\end{keyfactbox}

Three numbers (ATMF, 25\,RR, 25\,BF) pin three vol points: the ATMF
and 25-delta call/put vols.  Adding 10\,RR and 10\,BF gives five
points for interpolation.

\subsection{Example: EUR/USD 1-month}

With $\volat_{\mathrm{ATMF}} = 8.50\%, \mathrm{RR}_{25} = -0.40\%,
\mathrm{BF}_{25} = 0.18\%$:
\begin{align}
  \volat(25\Delta\text{-call}) &= 8.50 + 0.18 + \tfrac12(-0.40)
                                = 8.48\%, \\
  \volat(25\Delta\text{-put})  &= 8.50 + 0.18 - \tfrac12(-0.40)
                                = 8.88\%.
\end{align}
Negative 25\,RR means the put is more expensive than the call in
vol --- the market prices EUR downside (USD strength) as more
volatile than EUR upside.  Typical EUR/USD skew; opposite in USD/JPY
where JPY-strength tails command more vol.

\begin{intuitionbox}
ATMF = height, RR = tilt, BF = curvature above the straight line
through the wings --- the minimum three-parameter smile.  Traders
express directional vol views via $\Delta$RR (``buying 1M EUR/USD
risk-reversal'' = betting EUR-downside vol widens further over
EUR-upside) and convexity views via $\Delta$BF.
\end{intuitionbox}

\subsection{Smile reconstruction: five-point interpolation}
\label{subsec:smile_reconstruction}

For a more refined smile, desks quote a five-point system
$(\volat_{\mathrm{ATMF}}, \mathrm{RR}_{10}, \mathrm{BF}_{10},
\mathrm{RR}_{25}, \mathrm{BF}_{25})$.  A typical 3-month EUR/USD
quote block might read
\begin{center}
\small
\begin{tabular}{@{}ll@{}}
\toprule
ATMF & $9.00\%$ \\
25$\Delta$ RR & $-0.50\%$ \\
25$\Delta$ BF & $0.20\%$ \\
10$\Delta$ RR & $-1.10\%$ \\
10$\Delta$ BF & $0.75\%$ \\
\bottomrule
\end{tabular}
\end{center}
from which the five wing vols are
\begin{align}
  \volat(25\Delta\text{-call}) &= 9.00 + 0.20 - 0.25 = 8.95\%, \\
  \volat(25\Delta\text{-put})  &= 9.00 + 0.20 + 0.25 = 9.45\%, \\
  \volat(10\Delta\text{-call}) &= 9.00 + 0.75 - 0.55 = 9.20\%, \\
  \volat(10\Delta\text{-put})  &= 9.00 + 0.75 + 0.55 = 10.30\%.
\end{align}
These five points are interpolated into a $\volat(\Kstrike)$ curve.
Standard techniques: Vanna--Volga \citep{castagna2007consistent},
SABR \citep{hagan2002managing}, and monotone-convex splines for
no-arbitrage.  They agree on interpolation but diverge on
extrapolation beyond 10-delta; FX smiles are thus often quoted to
5-delta when deep-OTM options trade.

\begin{pitfallbox}[RR sign conventions disagree across regions]
London desks quote RR as call minus put; some NY desks quote put
minus call, flipping every historical series.  Check the sign on a
broker sheet against implied call/put vols before taking a view.
Sanity check: EUR/USD 25\,RR is almost always \emph{negative} on
London convention; a positive number without an inversion means
the sign is flipped.
\end{pitfallbox}


%% ================================================================
\section{FX market structure: CLS and dealer tiering}
\label{sec:fx_market_structure}
%% ================================================================

FX trades in an over-the-counter market split into three tiers.

\textbf{Tier 1} --- interbank: the G10 banks (JPM, Citi, DB, UBS,
Goldman, HSBC, BofA, BNP Paribas, and a few others) making two-way
prices in size to each other and to top clients.  Quotes trade on
EBS and Refinitiv FX Matching; inside spread on EUR/USD is sub-pip
for \$5M clips.

\textbf{Tier 2} --- single-/multi-dealer platforms: each bank's
e-FX portal (Barx, NeoBank, Autobahn, \ldots) and ECNs (FXall, 360T,
Hotspot, Currenex).  Asset managers, hedge funds, and corporates
trade RFQ or streaming, at spreads scaling with size, urgency, and
toxicity.

\textbf{Tier 3} --- retail: CFD brokers, spreads 5--20$\times$
interbank and often 2\% per trip on exotics.

\textbf{Settlement.}  Most FX cashflows settle through CLS
(Continuous Linked Settlement), launched 2002, settling 18
currencies PvP in a daily five-hour window.  CLS eliminated
``Herstatt risk'' --- your day-side payment going out before the
counterparty fails its other-side leg (a live risk before 2002;
Bankhaus Herstatt's 1974 \$620M failure is the namesake).  ITM
options settle via spot, so non-CLS pairs carry meaningfully higher
operational and credit risk.

\begin{historybox}[title={Herstatt and the birth of CLS}]
26 June 1974: German regulators closed Bankhaus Herstatt at 15:30
CET.  US banks had paid DEM earlier that morning for DEM/USD trades,
expecting USD via CHIPS later.  Herstatt's USD leg never paid,
leaving roughly \$620M of chain losses.  The episode drove the
three-decade project culminating in CLS.  Non-CLS pairs remain
the ones with highest implied vol and largest overnight risk
charges.
\end{historybox}

%% ================================================================
\section{Vanna--Volga adjustments}
\label{sec:vanna_volga}
%% ================================================================

The \textbf{Vanna--Volga} method \citep{castagna2007consistent} is
the standard FX technique for pricing exotics off the three
ATMF/25RR/25BF inputs.  It is a first-order smile correction: start
with the BS/GK price at ATMF vol and add terms making the three
vanilla hedges (ATMF straddle, 25$\Delta$ RR, 25$\Delta$ BF) match
market exactly.

Vanna and volga are the mixed and pure second derivatives in
$\volat$ and $\Spx$:
\begin{equation}
  \mathrm{vanna} \;=\;
    \frac{\partial^2 V}{\partial \Spx\, \partial \volat},
  \qquad
  \mathrm{volga} \;=\;
    \frac{\partial^2 V}{\partial \volat^2}.
\end{equation}
A risk-reversal is mostly vanna (long delta--vol correlation); a
butterfly is mostly volga (long vol-of-vol).  The VV price is the
BS/GK price plus a linear combination of vanna, volga, and vega with
coefficients pinned by matching the three hedging vanillas:
\begin{equation}
  V_{\mathrm{VV}} \;=\;
  V_{\mathrm{BS}} + w_{\mathrm{ATMF}}\,\mathrm{vega}
                   + w_{\mathrm{BF}}\,\mathrm{volga}
                   + w_{\mathrm{RR}}\,\mathrm{vanna}.
\end{equation}
The weights $w_\cdot$ solve a closed-form 3$\times$3 linear system
requiring that three reference vanillas (ATMF straddle, 25$\Delta$
call, 25$\Delta$ put) reproduce their market premia.

\begin{keyfactbox}[Vanna--Volga in one line]
\label{fact:vanna_volga}
Correct the Black--Scholes price of an exotic by adding the vanna
and volga dollar-risk, weighted by the market-price surplus over
BS of a 25$\Delta$ risk-reversal and butterfly respectively, plus
a vega correction for the ATMF skew of the exotic's own vol.
Three vanilla hedges, three corrections, exact match on those
three, first-order smile consistency everywhere else.
\end{keyfactbox}

VV is not rigorous (not a no-arbitrage price under any model) but
reproduces desk quotes for single-barrier options to within a basis
point in normal markets.  Multi-barrier and path-dependent exotics
use smile-consistent MC or PDE instead.


%% ================================================================
\section{Desk and competition context}
\label{sec:fx_context}
%% ================================================================

\subsection{Choosing the right FX hedge}

For a hedger with a specific view and exposure, the right FX
product follows from three questions: horizon, desired payoff shape,
tolerance for basis risk.  Short decision guide:
\begin{center}
\small
\begin{tabular}{@{}p{0.35\textwidth} p{0.60\textwidth}@{}}
\toprule
\textbf{Hedger's situation} & \textbf{Standard product} \\
\midrule
Certain receivable / payable, single date &
  Forward (if deliverable); NDF (if not) \\
Budget-rate floor, willing to cap upside &
  Collar (long put + short call at matching deltas) \\
Protection against extreme move only &
  Out-of-the-money put / call \\
``Cheap'' hedge, strong directional view &
  Knock-out option (accept tail risk on hedge disappearing) \\
Known range view (spot stays in $[a,b]$) &
  No-touch option (pays if range holds) \\
Series of monthly cashflows &
  Participating forward / accumulator / (beware) TARF \\
Equity-index exposure in different currency &
  Quanto option or FX-hedged share class \\
\bottomrule
\end{tabular}
\end{center}
All price off the same GK engine plus the smile of
\cref{sec:fx_delta_conventions}; only the payoff function changes.


\begin{prosperitybox}[title={Prosperity Round 1 manual: triangular FX}]
Round 1's manual gives a small cluster of rates between fictional
currencies and asks for a trade sequence ending with more shells
than you started with --- triangular FX arb
(\cref{ch:arbitrage_zoo}) in disguise.  A cycle
$\mathrm{A} \to \mathrm{B} \to \mathrm{C} \to \mathrm{A}$ is
profitable iff the product of rates exceeds 1 net of fees.  The
Frankfurt Hedgehogs (\cref{ch:top_team_lessons}) brute-forced
4-legged paths.  Real-world FX arb scales this to the full
cross-rate graph with transaction costs, latency, and basis.
\end{prosperitybox}

\begin{gsbox}[title={The FX options desk}]
On a sell-side FX options desk, the Strat owns three systems.
First, a \emph{vol surface builder} ingesting ATMF/RR/BF broker
quotes across tenors and delta slices --- usually a SABR or VV fit
\citep{castagna2007consistent}.  Second, a \emph{pricer} for
vanillas (GK) and barriers/digitals/touches (closed-form with
skew corrections).  Third, a \emph{MC exotics pricer} for TARFs,
accumulators, and path-dependent structures, wired into the firm's
risk system for aggregate Greeks and overnight VaR.  Interviews
routinely ask you to derive CIP, state GK, explain ATMF/RR/BF, and
sketch a knock-out pricer.
\end{gsbox}


%% ================================================================
\section{Key facts to memorise}
\label{sec:fx_summary}
%% ================================================================

\begin{itemize}
  \item \textbf{FX is the largest financial market}: \$7.5T per
    day (BIS 2022), larger than all equity markets combined.
  \item A pair $\mathrm{BASE}/\mathrm{QUOTE}$ is quoted in units of
    quote per unit of base.  EUR/USD = 1.08 means one EUR buys
    1.08 USD.
  \item A \textbf{pip} is $10^{-4}$ for most pairs, $10^{-2}$ for
    JPY pairs.  \textbf{Forward points} = $F - S$ in pips.
  \item \textbf{Covered interest parity}:
    $F = S\,e^{(r_d - r_f)T}$.  The high-yielding currency trades
    at a forward \emph{discount} (the low-yielder at a premium).
  \item Post-2008 CIP deviates by 10--50~bp persistently.  The
    wedge is the \emph{cross-currency basis}, traded via XCCY
    basis swaps.
  \item \textbf{Garman--Kohlhagen}: the FX Black--Scholes.  Treats
    foreign currency as a continuous-yield asset with yield $r_f$.
  \item \textbf{Put--call parity} in FX:
    $\Call - \Put = S e^{-r_f T} - \Kstrike e^{-r_d T}$.
  \item \textbf{Quanto adjustment}: subtract
    $\rho \volat_{\mathrm{asset}} \volat_{\mathrm{FX}}$ from the
    drift when the payoff is in a currency other than the
    underlying's natural one.
  \item \textbf{Barrier options}: knock-in + knock-out = vanilla.
    Cheaper than vanillas but disappear when the market moves far
    enough.
  \item \textbf{TARFs} cap client upside and leverage client
    downside --- the structural mis-selling trap that cost Asian
    corporates \$10bn+ in 2015.
  \item \textbf{FX vol is quoted in delta space}: ATMF vol, plus
    25-delta risk-reversal (skew) and 25-delta butterfly
    (convexity).  Three numbers pin down three points on the smile.
  \item \textbf{Negative 25\,RR} = market prices downside (for the
    base currency) at higher vol than upside.  Standard in
    EUR/USD; opposite in many USD/EM pairs.
\end{itemize}
